% =====================================================================
%  序言（原文 PDF p.9–14，罗马页码 v–x）
%  译者：Lead
% =====================================================================
\chapter*{序言}
\markboth{序言}{序言}
\addcontentsline{toc}{chapter}{序言}

Langlands 纲领近年来已成为数学“大统一理论”的蓝图。它最初是作为数论与自守表示之间的桥梁提出的，如今已扩展到几何与量子场论等领域，把看似无关的学科编织成一张充满诱人猜想的网络。Langlands 对应在这些数学与物理的不同领域中以多种方式显现，但那些相同的突出特征——例如 Langlands 对偶群的出现——始终存在。这指向某种深刻而难以捉摸的东西，而这正是该对应如此引人入胜的原因。

Langlands 纲领的一个常见主题是局部图景与整体图景之间的相互作用。例如，在数论语境中，“整体”指的是一个数域（有理数域的有限扩张）及其伽罗瓦群，而“局部”指的是一个局部域，例如 $p$-adic 数域，连同其伽罗瓦群。在 Langlands 对应的另一侧，整体情形下我们有自守表示，局部情形下我们有约化群（例如 $GL_n$）在局部域上的表示。

在几何语境中，角色阵容改变了：在伽罗瓦一侧，整体情形下我们现在有复黎曼面 $X$ 上带平坦联络的向量丛，局部情形下则是 $X$ 上一点周围的去心圆盘 $D^\times$ 上带平坦联络的向量丛。几何 Langlands 对应另一侧对象的定义则更为微妙。当丛上的平坦联络没有奇点时，这一侧已得到较好的理解（在 A. Beilinson、V. Drinfeld、G. Laumon 等人的工作之后）。此时相应的对象是 $X$ 上向量丛模空间上所谓的“Hecke 本征层”。它们是无分歧自守函数的几何类比。无分歧的整体几何 Langlands 对应应当把丛上的一个平坦联络（无奇点）映到一个 Hecke 本征层。（除其他文献外，最近的综述 [Frenkel 2007] 讨论了这一点，读者可参阅其中细节。）

然而，在更一般的带分歧联络的情形，即带奇点的情形，几何 Langlands 对应无论在局部还是整体情形都要神秘得多。实际上，推动力此时更多地转向局部情形。这是因为我们考虑的平坦联络在黎曼面上只有有限多个奇点。整体的带分歧对应在很大程度上由这些点周围去心圆盘上发生的事情决定，而那属于局部对应的范围。所以问题实际上变成：局部 Langlands 对应的几何类比是什么？

经典局部 Langlands 对应把 $p$-adic 群的表示与伽罗瓦表示联系起来。在几何版本中，我们应当把 $p$-adic 群替换为（形式）圈群 $G(\!(t)\!)$，即从（形式）去心圆盘 $D^\times$ 到复约化代数群 $G$ 的映射构成的群。伽罗瓦表示应当替换为 $D^\times$ 上带平坦联络的向量丛。这些就是局部几何 Langlands 参数。对其中每个参数，我们都应当能附上一个形式圈群的表示。

最近，Dennis Gaitsgory 与我提出了一个描述圈群的这些表示的一般方案。我们的方案中一个重要的新要素是：与经典对应不同，圈群现在作用在范畴上，而不是向量空间上。因此，圈群的 Langlands 对应是范畴性的：我们把 $G(\!(t)\!)$ 的范畴表示配给局部 Langlands 参数。我们已经提出如何利用仿射 Kac--Moody 代数 $\hat{\mathfrak{g}}$（它是圈代数 $\mathfrak{g}(\!(t)\!)$ 的中心扩张）的表示来构造这些范畴。因此，局部几何 Langlands 对应看来是 Langlands 哲学与仿射 Kac--Moody 代数表示论成功结合的产物。

仿射 Kac--Moody 代数有一个参数，称为水平。对该参数的一个特殊值——临界水平——仿射 Kac--Moody 代数的完备泛包络代数获得一个异常大的中心。1991 年，Boris Feigin 与我证明了这个中心典范同构于 $D^\times$ 上 opers 空间上的函数代数。opers 是 $D^\times$ 上带平坦联络的丛，并带一个附加数据（由 Drinfeld--Sokolov 与 Beilinson--Drinfeld 定义）。值得注意的是，它们的结构群不是 $G$，而是 Langlands 对偶群 ${}^L G$，这与 Langlands 哲学的一般精神一致。这是核心结果；它意味着同样的突出特征渗透在 $p$-adic 群表示论与圈群的（范畴）表示论两者之中。这一结果原为 V. Drinfeld 的猜想，并且在他与 A. Beilinson 借助 Hitchin 系统量子化研究整体几何 Langlands 对应的方法中起重要作用。中心与 opers 上函数之间的同构意味着临界水平 $\hat{\mathfrak{g}}$ 的表示范畴“生活”在 $D^\times$ 上 ${}^L G$-opers 的空间之上，而圈群 $G(\!(t)\!)$ 则“纤维式地”作用在这个范畴上。简言之，Gaitsgory 与我的方案是：这个范畴的“纤维”正是与 ${}^L G$-opers 背后的局部 Langlands 参数相对应的、我们所要寻找的 $G(\!(t)\!)$ 的范畴表示。这对局部与整体几何 Langlands 对应、以及对 $\hat{\mathfrak{g}}$ 的表示论都有许多非平凡的推论。我们希望对这些范畴的进一步研究将给我们带来关于 Langlands 对应之谜的新线索与洞见。\footnote{我们注意到，A. Beilinson 对局部几何 Langlands 对应还有另一个方案 [2006]，使用整水平（小于临界）的仿射 Kac--Moody 代数的表示。理解他的方案与我们的方案之间的联系将是饶有趣味的。}

本书的目标是对圈群的局部几何 Langlands 对应以及相关的仿射 Kac--Moody 代数表示论作系统而自足的介绍。它涵盖了过去二十年在这一领域的研究，部分基于我 2002 年与 2004 年在加州大学伯克利分校讲授的研究生课程。书中整个理论从零开始建立，所有必需的概念都加以定义，所有本质结果都沿途证明。我们引入 Weil--Deligne 群、Langlands 对偶群、仿射 Kac--Moody 代数、顶点代数、射流概形、opers、Miura opers、屏蔽算子等概念，并用详细的例子加以说明。特别地，许多结果首先在最简单的 $SL_2$ 情形中解释。读者几乎不需要超出标准大学代数的背景（最后一章可能例外）；我们甚至在附录中解释了一些标准概念，例如泛包络代数。

在开篇一章，我们对经典 Langlands 对应作教学式的概述，并给出有动机的、逐步走向几何设定的过渡。这引导我们研究仿射 Kac--Moody 代数，特别是完备泛包络代数的中心。随后我们非常详细地回顾一族仿射 Kac--Moody 代数表示的构造，称为 \emph{Wakimoto 模}。它们是 Feigin 与我在 20 世纪 80 年代末追随 M. Wakimoto 的工作而定义的。这些模为我们发展仿射代数表示论提供了有效工具。特别地，它们在我们证明中心谱与 opers 同构的过程中至关重要。对 Wakimoto 模的详细阐述与这一同构的证明构成本书的主要部分。这些结果使我们能够在临界水平仿射 Kac--Moody 代数的表示论与 opers 的几何之间建立深刻的联系。在最后一章，我们回顾 Gaitsgory 与我关于在 Langlands 对应框架中描述与 opers 相联系的表示范畴的结果与猜想。我还讨论这对整体几何 Langlands 对应的含义。这些只是一个新理论的最初几步，我们希望它最终能帮助我们揭示 Langlands 对偶的秘密。

% [Tutor Note: 序言为全书定调，核心线索只有一条：临界水平 k=-h^\vee 处 \hat{\mathfrak g} 的完备泛包络代数有一个“异常大”的中心，它同构于 D^\times 上 {}^L G-opers 的函数环（Feigin--Frenkel 定理）。读懂全书的关键是把“中心=opers 上函数”这一代数事实翻译成几何语言：临界水平的表示范畴纤维化在 opers 空间上，而 G((t)) 沿纤维作用。第 8 章完成该同构的证明，第 10 章把它升级为范畴化的 Langlands 对应。]

\noindent\textbf{内容提要。}以下更详细地介绍本书内容。

第 1 章是本书主题的导引。我们首先概述经典设定下的局部与整体 Langlands 对应。由于整体情形已在我的近期综述 [Frenkel 2007] 中详细讨论，这里我主要集中于局部情形。接着，我解释为了把局部 Langlands 对应搬到几何与圈群表示论的领域需要作哪些改变。我教学式地讲述伽罗瓦群、带联络的主丛与中心扩张等主题。这一讨论把我们引向如下问题：如何给每个局部几何 Langlands 参数配上一个带形式圈群作用的阿贝尔范畴？

在第 2 章，我们在仿射 Kac--Moody 代数表示论的语境中处理这个问题。这促使我们研究 $\hat{\mathfrak{g}}$ 的完备泛包络代数的中心。我们首先用初等方法来做，但很快意识到需要更精细的技术。这一技术就是顶点代数理论。我们给出顶点代数的速成课程（遵循 [Frenkel and Ben-Zvi 2004]），总结所有必要的概念与结果。

有了这些结果，我们在第 3 章开始更深入地研究临界水平 $\hat{\mathfrak{g}}$ 的完备泛包络代数的中心（我们会发现，离开临界水平中心是平凡的）。我们描述最简单仿射 Kac--Moody 代数 $\hat{\mathfrak{sl}}_2$ 情形的中心，以及任意 $\hat{\mathfrak{g}}$ 的中心的准经典类比。

在第 4 章，我们引入 oper 这一关键的几何概念；它最初出现于 [Drinfeld and Sokolov 1985; Beilinson and Drinfeld 1997]。我们陈述主要结果（属于 Feigin 与我 [1992; 2005b]）：与 $\hat{\mathfrak{g}}$ 相应的临界水平中心同构于 ${}^L G$-opers 上的函数代数。为证明这一结果，我们需要发展 Wakimoto 模理论。这将在第 5、6 章完成，遵循 [Frenkel 2005b]。我们首先解释有限维单李代数的类似理论，它作为我们构造的原型。然后解释无限维情形中的非平凡要素，例如在圈空间上的微分算子代数中实现圈代数时所遇到的上同调障碍。这导致对临界水平非平凡性的概念性解释。在第 6 章，我们完成 Wakimoto 模的构造，包括临界与非临界水平。我们证明关于仿射 Kac--Moody 代数表示的一些有用结果，例如 Kac--Kazhdan 猜想。

在建立了 Wakimoto 模理论之后，我们准备处理中心与 opers 上函数代数之间的同构。在第 7 章开头，我们给出该同构证明的详细概览。第 7 章其余部分引入 Wakimoto 模之间一类重要的交织算子，称为屏蔽算子。我们在第 8 章利用这些算子以及关于分次伴随代数的一些结果来完成主要结果的证明，并把中心等同于 opers 上的函数（这里遵循 [Frenkel 2005b]）。特别地，我们澄清 Langlands 对偶群在这一同构中出现的原因，把它追溯到顶点代数之间一种称为 W-代数的对偶。在本章末尾，我们讨论中心上的顶点 Poisson 结构，并把中心在 Wakimoto 模上的作用等同于 Miura 变换。

在第 9 章，我们更深入地研究临界水平仿射 Kac--Moody 代数的表示。我们首先（遵循 [Beilinson and Drinfeld 1997] 与 [Frenkel and Gaitsgory 2006c]）引入 $D^\times$ 上 opers 空间的某些子空间：带正则奇点的 opers 与幂零 opers，并解释它们之间的相互关系。然后（遵循 [Frenkel and Gaitsgory 2006c]）讨论带正则奇点的 Miura opers 以及 Miura 变换在它们上面的作用。最后，我们描述 [Frenkel and Gaitsgory 2006c] 与 [Frenkel and Gaitsgory 2007c] 中关于临界水平 Verma 模与 Weyl 模的自同态代数的结果。

在第 10 章，我们把本书的结果汇集起来，解释 Gaitsgory 与我提出的局部几何 Langlands 对应方案。我们回顾我们合作诸文 [Frenkel and Gaitsgory 2004; 2006c; 2006a; 2005; 2006b; 2007c] 中的结果与猜想，强调几何 Langlands 对应与经典 Langlands 对应之间的类比。我们详细讨论 opers 与局部系之间的相互作用。然后考虑最简单的局部系，即平凡局部系。相应的 $G(\!(t)\!)$ 的范畴表示是 $p$-adic 群无分歧表示的类比。即使在这一情形，我们也会看到相当非平凡的要素，它们模仿经典理论中的相应要素，同时以非平凡的方式推广它们。下一个例子要复杂得多：$D^\times$ 上带正则奇点与幂幺单值性的局部系。它们是 $p$-adic 域伽罗瓦群的温和分歧表示的类比。相应的范畴结果与 Springer 纤维上的拟凝聚层范畴密切相关，而 Springer 纤维是 Langlands 对偶群旗簇的代数子簇。我们总结 [Frenkel and Gaitsgory 2006c] 关于这些范畴的猜想与结果，并用 $\hat{\mathfrak{sl}}_2$ 情形的显式计算加以说明。我们还提出这一方向上的一些未解决问题。最后，我们讨论这一方法对整体 Langlands 对应的含义。

\noindent\textbf{致谢。}我感谢 Boris Feigin 与 Dennis Gaitsgory 在本书所综述的我们合作工作中的协作。

多年来，我在与 Alexander Beilinson、David Ben-Zvi、Joseph Bernstein、Vladimir Drinfeld、Victor Kac 和 David Kazhdan 关于 Langlands 对应与仿射 Kac--Moody 代数表示的无数讨论中获益良多。我对他们所有人表示感谢。我还要感谢 R. Bezrukavnikov 与 V. Ginzburg 对最后一章材料的某些问题提出的有用意见。

我感谢 Alex Barnard 允许我使用他在我于加州大学伯克利分校开设的研究生课程上所作的笔记，来准备第 2、3 章。我还感谢 Arthur Greenspoon、Reimundo Heluani、Valerio Toledano Laredo 与 Xinwen Zhu 对本书草稿提出的宝贵意见。我感谢剑桥大学出版社的 David Tranah 在印刷版出版过程中给予的专业指导与友好建议。我感谢 Mathematical Sciences Publishers，特别是 Paulo Ney de Souza，为我的书制作了出色的、带完整超链接的电子版。

在撰写本书过程中，我得到了 DARPA 通过其“理论数学的基础性进展”计划的慷慨资助。我还部分得到美国国家科学基金会 DMS-0303529 号项目的资助。我对此深表感谢。
